% Part: first-order-logic
% Chapter: completeness
% Section: identity

\documentclass[../../../include/open-logic-section]{subfiles}

\begin{document}

\olfileid{fol}{com}{ide}
\olsection{Identiteit}

\begin{explain}
Met het termmodel uit de vorige paragraaf kunnen we de volledigheid van
de predicatenlogica bewijzen voor verzamelingen~$\Gamma$ zonder~$\eq$.
Dat model vervult elke $!A \in \Gamma^*$ zonder~$\eq$, en dus ook elke
$!A \in \Gamma$. Zodra~$\eq$ wel voorkomt, ontstaat een probleem.
De verzameling $\Gamma^*$ kan dan !!a{sentence}~$\eq[t][t']$ bevatten.
Maar in het termmodel is de waarde van een term gewoon die term zelf.
Zijn $t$ en $t'$ verschillende termen, dan zijn hun waarden dus ook
verschillend: respectievelijk $t$ en $t'$. Daardoor is $\eq[t][t']$
onwaar. We lossen dit op met ``quotiëntvorming'': termen die volgens de
verzameling gelijk zijn, behandelen we als één object.
\end{explain}

\begin{defn}
  Neem een consistente en !!{complete} verzameling zinnen $\Gamma^*$ in
  $\Lang L$. We definiëren de relatie $\approx$ op de gesloten termen
  van~$\Lang L$ door
  \[
  t \approx t' \text{\quad precies als \quad} \eq[t][t'] \in \Gamma^*
  \]
\end{defn}

\begin{prop}
\ollabel{prop:approx-equiv}
De relatie $\approx$ heeft deze eigenschappen:
\begin{enumerate}
\item $\approx$ is reflexief.
\item $\approx$ is symmetrisch.
\item  $\approx$ is transitief.
\item Als $t \approx t'$, als $f$ !!a{function} is en als $t_1$, \dots,
  $t_{i-1}$, $t_{i+1}$, \dots, $t_n$ gesloten termen zijn, dan
\[
\Atom{f}{t_1,\dots, t_{i-1}, t, t_{i+1}, \dots, t_n} \approx
\Atom{f}{t_1,\dots, t_{i-1}, t', t_{i+1}, \dots, t_n}.
\]
\item Als $t \approx t'$, als $R$ !!a{predicate} is en als $t_1$, \dots,
  $t_{i-1}$, $t_{i+1}$, \dots, $t_n$ gesloten termen zijn, dan
\begin{multline*}
\Atom{R}{t_1,\dots, t_{i-1}, t, t_{i+1}, \dots, t_n} \in \Gamma^* \text{ precies als } \\
\Atom{R}{t_1,\dots, t_{i-1}, t', t_{i+1}, \dots, t_n} \in \Gamma^*.
\end{multline*}
\end{enumerate}
\end{prop}

\begin{proof}
De verzameling $\Gamma^*$ is consistent en !!{complete}. Daarom geldt
$\eq[t][t'] \in \Gamma^*$ precies als
$\Gamma^* \Proves \eq[t][t']$. We hoeven dus alleen het volgende te bewijzen:
\begin{enumerate}
\item $\Gamma^* \Proves \eq[t][t]$ voor elke gesloten term~$t$.
\item Als $\Gamma^* \Proves \eq[t][t']$, dan $\Gamma^* \Proves \eq[t'][t]$.
\item Als $\Gamma^* \Proves \eq[t][t']$ en $\Gamma^* \Proves
  \eq[t'][t'']$, dan $\Gamma^* \Proves \eq[t][t'']$.
\item Als $\Gamma^* \Proves \eq[t][t']$, dan
\[
\Gamma^* \Proves
\eq[\Atom{f}{t_1,\dots,t_{i-1},t,t_{i+1},,\dots,t_n}][\Atom{f}{t_1,\dots,t_{i-1},t',t_{i+1},\dots,t_n}]
\]
voor elke $n$-plaatsige !!{function}~$f$ en gesloten termen $t_1$, \dots,
$t_{i-1}$, $t_{i+1}$, \dots,~$t_n$.
\item Als $\Gamma^* \Proves \eq[t][t']$ en
$\Gamma^* \Proves
\Atom{R}{t_1,\dots,t_{i-1},t,t_{i+1},\dots,t_n}$, then
$\Gamma^* \Proves \Atom{R}{t_1,\dots,t_{i-1},t',t_{i+1},\dots,t_n}$
voor elk $n$-plaatsig !!{predicate}~$R$ en gesloten termen $t_1$, \dots,
$t_{i-1}$, $t_{i+1}$, \dots,~$t_n$.
\end{enumerate}
\end{proof}

\begin{prob}
Maak het bewijs van~\olref[fol][com][ide]{prop:approx-equiv} af.
\end{prob}

\begin{defn}
Neem een consistente en !!{complete} verzameling $\Gamma^*$ in een
taal~$\Lang L$, een gesloten term $t$ en de relatie $\approx$ uit de
vorige definitie. Dan:
\[
\equivrep{t}{\approx} = \Setabs{t'}{t'\in \Trm[L], t \approx t'}
\]
and $\equivclass{\Trm[L]}{\approx} = \Setabs{\equivrep{t}{\approx}}{t \in \Trm[L]}$.
\end{defn}

\begin{defn}
\ollabel{defn:term-model-factor}
Neem het termmodel $\Struct M = \Struct M(\Gamma^*)$ voor~$\Gamma^*$ uit
\olref[mod]{defn:termmodel}. De !!{structure}
$\Struct{\equivclass{M}{\approx}}$ wordt dan als volgt opgebouwd:
\begin{enumerate}
\item $\Domain{\equivclass{M}{\approx}} = \equivclass{\Trm[L]}{\approx}$.
\item $\Assign{c}{\equivclass{M}{\approx}} = \equivrep{c}{\approx}$
\item $\Assign{f}{\equivclass{M}{\approx}}(\equivrep{t_1}{\approx}, \dots,
  \equivrep{t_n}{\approx}) = \equivrep{\Atom{f}{t_1,\dots, t_n}}{\approx}$
\item $\tuple{\equivrep{t_1}{\approx}, \dots, \equivrep{t_n}{\approx}} \in
  \Assign{R}{\equivclass{M}{\approx}}$ precies als
  $\Sat{M}{\Atom{R}{t_1,\dots, t_n}}$, en dus precies als
  $\Atom{R}{t_1,\dots,
  t_n} \in \Gamma^*$.
\end{enumerate}
\end{defn}

\begin{explain}
Bij de definities van $\Assign{f}{\equivclass{M}{\approx}}$ en
$\Assign{R}{\equivclass{M}{\approx}}$ schrijven we de elementen van
$\equivclass{\Trm[L]}{\approx}$ als $\equivrep{t}{\approx}$. We kiezen
dus een \emph{representant}~$t \in \equivrep{t}{\approx}$. De uitkomst
mag niet afhangen van die keuze. Kiezen we een andere term~$t'$ uit
dezelfde equivalentieklasse
($\equivrep{t}{\approx} = \equivrep{t'}{\approx}$), dan moeten de
definities hetzelfde resultaat geven. Neem bijvoorbeeld een eenplaatsig
!!{predicate}~$R$. De laatste regel van de definitie zegt dat
$\equivrep{t}{\approx} \in \Assign{R}{\equivclass{M}{\approx}}$ precies
als $\Sat{M}{\Atom{R}{t}}$. Stel nu dat voor een andere term~$t'$ met
$t \approx t'$ juist $\Sat/{M}{\Atom{R}{t}}$. Dan zou de definitie eisen
dat $\equivrep{t'}{\approx} \notin
\Assign{R}{\equivclass{M}{\approx}}$. Maar uit $t \approx t'$ volgt
$\equivrep{t}{\approx} = \equivrep{t'}{\approx}$. Hetzelfde element zou
dan zowel $\equivrep{t}{\approx} \in
\Assign{R}{\equivclass{M}{\approx}}$ als $\equivrep{t}{\approx} \notin
\Assign{R}{\equivclass{M}{\approx}}$ moeten opleveren. Dat kan niet.
Volgens \olref{prop:approx-equiv} ontstaat deze situatie ook werkelijk
niet.
\end{explain}

\begin{prop}
$\Struct{\equivclass{M}{\approx}}$ is goed gedefinieerd. Dat betekent het
volgende. Zijn $t_1$, \dots, $t_n$, $t_1'$, \dots, $t_n'$ gesloten termen
en geldt $t_i \approx t_i'$, dan
\begin{enumerate}
\item $\equivrep{\Atom{f}{t_1,\dots, t_n}}{\approx} =
    \equivrep{\Atom{f}{t_1',\dots, t_n'}}{\approx}$, dus
  \[
  \Atom{f}{t_1,\dots, t_n} \approx \Atom{f}{t_1',\dots, t_n'}
  \]
  en
\item $\Sat{M}{\Atom{R}{t_1,\dots, t_n}}$ precies als
  $\Sat{M}{\Atom{R}{t_1',\dots, t_n'}}$, dus
  \[
    \Atom{R}{t_1,\dots, t_n} \in \Gamma^* \text{ precies als }
    \Atom{R}{t_1',\dots, t_n'} \in \Gamma^*.
  \]
\end{enumerate}
\end{prop}

\begin{proof}
Dit volgt uit \olref{prop:approx-equiv} met inductie naar~$n$.
\end{proof}

Net als bij het termmodel hebben we eerst het volgende lemma nodig. Daarna
kunnen we het waarheidslemma bewijzen.

\begin{lem} 
\ollabel{lem:val-in-termmodel-factored} Neem $\Struct M = \Struct M
 (\Gamma^*)$. Dan geldt $\Value{t}{\equivclass{M}{\approx}} = \equivrep{t}
 {\approx}$.
\end{lem}
\begin{proof} 
Dit bewijs werkt net als dat van \olref[mod]{lem:val-in-termmodel}.
\end{proof}

\begin{prob}
Maak het bewijs van~\olref[fol][com][ide]{lem:val-in-termmodel-factored} af.
\end{prob}

\begin{lem}
\ollabel{lem:truth}
$\Sat{\equivclass{M}{\approx}}{!A}$ precies als $!A \in \Gamma^*$, voor
elke zin~$!A$.
\end{lem}

\begin{proof}
We gebruiken inductie naar~$!A$, net als in het bewijs van
\olref[mod]{lem:truth}. Alleen het geval $!A \ident \eq[t][t']$ vraagt
extra aandacht.
\begin{align*}
\Sat{\equivclass{M}{\approx}}{\eq[t][t']} & \text{ precies als } \equivrep{t}{\approx} = \equivrep{t'}{\approx}
\text{ (volgens de definitie van $\Struct{\equivclass{M}{\approx}}$)}\\
& \text{ precies als } t \approx t' \text{ (volgens de definitie van $\equivrep{t}{\approx}$)}\\
& \text{ precies als } \eq[t][t'] \in \Gamma^* \text{ (volgens de definitie van $\approx$).}
\end{align*}
\end{proof}

\begin{digress}
Het model $\Struct{M(\Gamma^*)}$ is altijd !!{enumerable} en oneindig.
Het quotiëntmodel $\Struct{\equivclass{M}{\approx}}$ kan toch eindig zijn,
want er kunnen maar eindig veel klassen $\equivrep{t}{\approx}$ zijn.
Dat is niet vreemd. De verzameling $\Gamma$ kan !!{sentence}s bevatten
die alleen waar kunnen zijn in een eindige !!{structure}. Een voorbeeld
is $\lforall[x][\lforall[y][x = y]]$. Deze !!{sentence} is consistent,
maar wordt alleen vervuld in !!{structure}s waarvan het !!{domain}
precies één !!{element} bevat.
\end{digress}

\end{document}
